\documentclass[10pt,a4paper]{article}
\usepackage{rpgcover}
\usepackage{illustrated-panel}
\usepackage{tocloft}
\babelprovide[import]{ukrainian}
\setmainfont{FreeSerif}
\RPGSetRunningTitle{Останній кидок}
\RPGSetCoverTitleFont{FreeSerif}
\RPGSetCoverSubtitleFont{FreeSerif}
\RPGSetCoverTextFont{FreeSerif}
\setlength{\cftsecnumwidth}{4.5em}
\setlength{\cftsecindent}{0em}
\renewcommand{\cftsecfont}{\normalsize\bfseries}
\renewcommand{\cftsecpagefont}{\normalsize\bfseries}
\setlength{\cftbeforesecskip}{2pt}
\setlength{\parskip}{2.5pt plus .5pt minus .5pt}
\linespread{1.06}\selectfont
\emergencystretch=1em
\tolerance=1000
\hypersetup{pdftitle={Останній кидок — українське видання},pdfauthor={Jacopo Biscella}}
\newcommand{\RPGUkrainianContents}{%
\thispagestyle{rpgfrontmatter}\vspace*{4mm}
\begin{center}{\fontsize{26}{30}\selectfont\bfseries ЗМІСТ\par}\vspace{9pt}
{\color{rpgline}\rule{.76\linewidth}{.55pt}}\end{center}\vspace{4mm}
\begingroup\csname @starttoc\endcsname{toc}\endgroup
\vspace{5mm}
\begin{center}\includegraphics[width=.51\textwidth,height=45mm,keepaspectratio]{assets/last-throw-crosshatch-redrawn.png}\end{center}
\vfill
\begin{center}{\small\itshape Гра Якопо Бішелли\par}\vspace{5pt}
{\small Цей твір поширюється за ліцензією Creative Commons Attribution—ShareAlike 4.0 International (CC BY-SA 4.0). Ви можете поширювати та адаптувати його, зокрема з комерційною метою, за умови зазначення автора та поширення похідних творів на тих самих умовах.\par}\end{center}\clearpage}
\begin{document}
\selectlanguage{ukrainian}
\frenchspacing
\RPGStartFrontMatter
\RPGSetCoverWordSpace{.75em}
\RPGFrontCover{ОСТАННІЙ КИДОК}{ВИБІР, ДОЛЯ, НАСЛІДКИ}{Jacopo Biscella}{assets/last-throw-crosshatch-redrawn.png}
\RPGStartContents
\RPGUkrainianContents
\RPGStartMainMatter
\RPGChapter{Кубик}
Коли Протагоніст намагається виконати дію з непевним результатом, він кидає шестигранні кубики (d6). Кількість кубиків і спосіб визначення результату залежать від рівня відповідної Характеристики.
\section*{Рівні та кидки}
\begin{rpgtable}{Y|Y|Y}
\rpghead{Рівень} & \rpghead{Кидок} & \rpghead{Значення} \\ \hline
Недосвідчений & 2d6, взяти менший & Не ваша стихія \\ \hline
Компетентний & 1d6 & Базове володіння \\ \hline
Експерт & 2d6, взяти більший & Майстер своєї справи \\ \hline
\end{rpgtable}
Відповідне \textbf{Вміння} покращує кидок на один щабель. \textbf{Ситуативний штраф} (рани, темрява, виснаження) погіршує його на один щабель. Ці впливи взаємно скасовуються.
\begin{rpgtable}{Y|Y}
\rpghead{Ситуація} & \rpghead{Ефективний кидок} \\ \hline
Недосвідчений + штраф & Погано: 3d6, взяти менший \\ \hline
Недосвідчений & 2d6, взяти менший \\ \hline
Недосвідчений + вміння / Компетентний & 1d6 \\ \hline
Компетентний + вміння / Експерт & 2d6, взяти більший \\ \hline
Експерт + вміння & Винятково: 3d6, взяти більший \\ \hline
\end{rpgtable}
\section*{Цільове значення}
Результат порівнюють із \textbf{Цільовим значенням (ЦЗ)} — мінімальним числом, потрібним для успіху. Оповідач визначає ЦЗ відповідно до складності дії.
\begin{rpgtable}{Y|Y}
\rpghead{ЦЗ} & \rpghead{Складність} \\ \hline
3 & Проста \\ \hline
4 & Звичайна \\ \hline
5 & Складна \\ \hline
6 & Екстремальна \\ \hline
Критичний успіх & Диво — потрібні дві шістки \\ \hline
\end{rpgtable}
Оповідач вільно інтерпретує ступінь успіху або невдачі: 6 проти ЦЗ 4 — блискучий результат, а 4 проти ЦЗ 4 — успіх на межі.
\section*{Критичний успіх}
Якщо на кубиках випало \textbf{дві або більше шісток}, це \textbf{Критичний успіх} — надзвичайний результат. Для нього потрібно щонайменше два кубики, отже він можливий лише на рівні Експерта або вище.

Деякі дії — вибити голими руками залізну браму, зламати стародавні чари чи здійснити диво — \textbf{вимагають Критичного успіху}. Спробувати таку дію може лише персонаж із відповідним Вмінням.

\textbf{Правило виснаження:} після невдалої спроби дії, що вимагає Критичного успіху, її не можна повторювати, доки обставини суттєво не зміняться: відпочинок, новий підхід, допомога чи зміна ситуації. Що вважати достатньою зміною, вирішує Оповідач.
\section*{Імовірності з першого погляду}
\begin{rpgtable}{Y|Y|Y|Y}
\rpghead{Рівень} & \rpghead{Шанс ≥4} & \rpghead{Шанс 6} & \rpghead{Критичний} \\ \hline
Поганий (3d6↓) & ≈12\% & ≈0,5\% & — \\ \hline
Недосвідчений (2d6↓) & ≈25\% & ≈3\% & — \\ \hline
Компетентний (1d6) & ≈50\% & ≈17\% & — \\ \hline
Експерт (2d6↑) & ≈75\% & ≈31\% & ≈2,8\% \\ \hline
Винятковий (3d6↑) & ≈87\% & ≈42\% & ≈7\% \\ \hline
\end{rpgtable}
\RPGChapter{Створення протагоніста}
Кожного Протагоніста визначають \textbf{Раса}, \textbf{Клас}, чотири \textbf{Характеристики}, набір \textbf{Вмінь}, \textbf{Риса}, \textbf{Тягар} і \textbf{Очки здоров'я}.

Гравець має \textbf{10 очок}, які розподіляє між Очками здоров'я, Характеристиками та Вміннями. Клас визначає вартість, обмеження й доступні можливості.
\section*{Вартість характеристик}
\begin{rpgtable}{Y|Y|Y}
\rpghead{Підвищення} & \rpghead{Основна} & \rpghead{Неосновна} \\ \hline
Недосвідчений → Компетентний & 1 очко & 2 очки \\ \hline
Компетентний → Експерт & 2 очки & 4 очки \\ \hline
Експерт → Винятковий & 4 очки & 8 очок \\ \hline
\end{rpgtable}
Кожне Очко здоров'я коштує 1 очко. Кожне Вміння або Гримуар — також 1 очко.

Расовий бонус надає знижку 1 очко на придбання рівня Компетентний для Характеристики расової родини (безкоштовно, якщо вона збігається з основною характеристикою Класу). Знижка діє лише на цей перехід.
\section*{Риса}
Фраза, що передає сутність героя. Не має механічного ефекту, але дає Оповідачеві матеріал для історії та спрямовує гравця.
\begin{rpgnote}«Я нікого не залишаю позаду» · «Владу беруть, а не просять» · «Краще незручна правда, ніж добра брехня»\end{rpgnote}
\section*{Тягар}
Щось, що герой несе із собою: докори сумління, борг, таємниця чи прокляття. Оповідач використовує це для створення дилем.
\begin{rpgnote}«Я вбив невинного під час війни» · «Я завдячую життям тому, кого ненавиджу» · «Хтось носить моє обличчя»\end{rpgnote}
\RPGChapter{Характеристики та вміння}
Кожен персонаж має чотири Характеристики, які спочатку перебувають на рівні Недосвідчений. Гравець підвищує їх, витрачаючи очки.
\begin{rpgtable}{Y|Y}
\rpghead{Характеристика} & \rpghead{За що відповідає} \\ \hline
Сила & Фізична міць, витривалість, ближній бій \\ \hline
Спритність & Рухливість, реакція, потайність, прицілювання, ухиляння \\ \hline
Знання & Мислення, культура, знання, соціальний вплив \\ \hline
Диво & Зв'язок із надприродним, чудеса \\ \hline
\end{rpgtable}
Відповідне Вміння покращує кидок на один щабель. Будь-яке Вміння можна придбати, якщо відповідна Характеристика досягла щонайменше рівня Компетентний. Вартість розвитку Характеристики природно обмежує вибір.
\section*{Вміння Сили}
\textbf{Важка зброя} — дворучна зброя, сокири, бойові молоти.

\textbf{Щити} — активний захист щитом, а також використання його як зброї.

\textbf{Боротьба} — беззбройний бій, захвати й утримання.

\textbf{Руйнування} — знищення предметів, дверей, стін і барикад.

\textbf{Витривалість} — опір болю, втомі, отрутам і суворим умовам.
\section*{Вміння Спритності}
\textbf{Потайність} — рух непомітно для інших.

\textbf{Акробатика} — рівновага, стрибки, складні рухи.

\textbf{Злом замків} — відкривання замків і механізмів.

\textbf{Пастки} — знаходження й знешкодження небезпечних пристроїв.

\textbf{Влучність} — точне ураження віддалених цілей.

\textbf{Кишенькові крадіжки} — непомітне викрадення речей.

\textbf{Верхова їзда} — керування верховою твариною в русі та бою.

\textbf{Виживання} — орієнтування за інстинктом, полювання, вистежування, таборування.
\section*{Вміння Знання}
\textbf{Дипломатія} — переговори, компроміси, порозуміння сторін.

\textbf{Риторика} — переконання, натхнення, промови перед натовпом.

\textbf{Обман} — маніпуляції та брехня.

\textbf{Залякування} — соціальний і психологічний тиск.

\textbf{Торгівля} — торги, оцінювання майна.

\textbf{Медицина} — діагностика, лікування, стабілізація.

\textbf{Мови} — лінгвістика, давні тексти, спілкування.

\textbf{Історія} — минуле, символи, контексти.

\textbf{Право} — закони, угоди, бюрократія.

\textbf{Етикет} — правила дворів і звичаї.

\textbf{Ремесла} — виготовлення, ремонт, оцінювання речей.

\textbf{Орієнтування} — карти, зорі, планування шляхів.

\textbf{Травництво} — природні ліки, настої, мазі.

\textbf{Тактика} — військова стратегія, координація, планування.
\section*{Вміння Дива}
\textbf{Спрямування} — зосередження та вивільнення стихійної енергії.

\textbf{Ритуали} — обрядова магія, викликання, освячення, пакти.

\textbf{Алхімія} — перетворення речовин, еліксири, алхімічний вогонь.

\textbf{Окультизм} — некромантія, демони, прокляття, заборонене знання.

\textbf{Магічне чуття} — відчуття магії, пророцтва, погляд за завісу.

Вміння Дива також забезпечують використання Гримуарів: кожному Гримуару потрібне певне Вміння Дива. Крім цих Вмінь, Містик має доступ до Гримуарів, описаних в окремому розділі.
\RPGChapter{Класи}
Кожен Протагоніст належить до одного з чотирьох Класів. Клас визначає основну Характеристику, діапазон Очок здоров'я та доступні Вміння.
\begin{rpgtable}{Y|Y|Y|Y|Y}
\rpghead{Клас} & \rpghead{Основа} & \rpghead{ОЗ мін.} & \rpghead{ОЗ макс.} & \rpghead{Вміння} \\ \hline
Боєць & Сила & 4 & 7 & 5 (Сила) \\ \hline
Авантюрист & Спритність & 3 & 5 & 8 (Спритність) \\ \hline
Учений & Знання & 2 & 5 & 14 (Знання) \\ \hline
Містик & Диво & 2 & 4 & 5 (Диво) + Гримуари \\ \hline
\end{rpgtable}
\section*{Боєць}
Опора загону. Багато Очок здоров'я, винятковий доступ до важких обладунків і нищівної зброї. Список Вмінь короткий, але не обов'язковий: його перевага — вистояти й завдати удару.
\section*{Авантюрист}
Гнучкий і пристосовний. Середній запас здоров'я, широкий набір Вмінь. Розв'язує проблеми завдяки фізичній винахідливості: відмикає двері, ховається, влучно б'є, уникає пасток.
\section*{Учений}
Знання — це сила. Найдовший список Вмінь дозволяє стати лікарем, дипломатом, стратегом, купцем або шахраєм. У бою допомагає Медицині, Тактиці, Залякуванню та Риториці.
\section*{Містик}
Крихкий, але здатний на неможливе. Мало Очок здоров'я, зате є Гримуари — закляття, яких не можуть використовувати інші Класи. Кожен Гримуар — ризик: невдала магія вражає і Містика, і його союзників.
\section*{Приклади створення}
\begin{rpgnote}
\textbf{Боєць-дворф} — Сила Компетентний: 0 (расова знижка), потім Експерт: 2. 6 ОЗ: 6. Два Вміння: 2. Разом: 10.

\textbf{Авантюрист-ельф} — Спритність Компетентний: 0, потім Експерт: 2. 4 ОЗ: 4. Чотири Вміння: 4. Разом: 10.

\textbf{Учений-людина} — Знання Компетентний: 0, потім Експерт: 2. 3 ОЗ: 3. П'ять Вмінь: 5. Разом: 10.

\textbf{Містик-фея} — Диво Компетентний: 0, потім Експерт: 2. 2 ОЗ: 2. Три Вміння: 3. Три Гримуари: 3. Разом: 10.
\end{rpgnote}
\RPGEngravedPanel{Приклад створення — Каель, авантюрист-ельф}{assets/panel-08-engraved.png}{Лука обирає ельфа (раса Спритності), Авантюриста. Завдяки расовому бонусу Спритність одразу Компетентна безкоштовно. Він має 10 очок.\par\smallskip Спритність Експерт: 2 очки. 4 ОЗ: 4 очки. Вміння — Потайність, Злом замків, Пастки, Влучність: 4 очки. Разом: 10.\par\smallskip Риса: «Немає дверей, яких я не відчиню».\par Тягар: «Я вкрав те, чого не мав торкатися».}
\RPGChapter{Раси}
Раси об'єднані в чотири родини, кожна з яких пов'язана з Характеристикою. Раса дає знижку 1 очко на придбання рівня Компетентний для цієї Характеристики. Приклади необов'язкові: Оповідач і гравці самі вирішують, хто існує у світі, або створюють нові раси.
\begin{rpgtable}{Y|Y|Y}
\rpghead{Родина} & \rpghead{Характеристика} & \rpghead{Приклади} \\ \hline
Раси Сили & Сила & Дворфи, орки, мінотаври \\ \hline
Раси Спритності & Спритність & Ельфи, напіврослики, гобліни \\ \hline
Раси Знання & Знання & Люди, мінливці, сирени \\ \hline
Раси Дива & Диво & Феї, лепрекони, демони, ангели \\ \hline
\end{rpgtable}
Єдине механічне правило — родина: вона визначає, яка Характеристика отримує знижку. Дворф-містик дешевше розвиває Силу, але платить повну ціну за Диво.
\RPGEngravedPanel{Дворф-містик?}{assets/panel-09-engraved.png}{Ґрунтар — дворф із незвичним покликанням. Кремезний, бородатий, з руками коваля, але він відчуває магію в залізі. Раса дає йому Компетентну Силу безкоштовно, проте серце лежить до Дива. Воно коштує дорого: 2 очки до Компетентного, 4 — до Експерта. Та коли він спрямовує полум'я крізь свій молот, ніхто вже не сміється.}
\RPGChapter{Бій}
Жодних жорстких ходів, ініціативи чи сітки. Бій — це структурована розмова.
\section*{Перебіг}
\begin{enumerate}[leftmargin=*,itemsep=2pt]
\item Оповідач описує ситуацію й запитує: «Що ви робите?»
\item Кожен Протагоніст оголошує свою дію.
\item Кидають відповідну Характеристику проти ЦЗ, призначеного Оповідачем.
\item Оповідач описує наслідки, і світ реагує.
\item Повторюють до завершення конфлікту — перемоги, втечі, капітуляції чи смерті.
\end{enumerate}
Немає ініціативи й фіксованого порядку. Оповідач надає слово відповідно до логіки сцени: хто ближчий, кому загрожує небезпека, чия дія найтерміновіша. У хаосі всі можуть оголосити дії одночасно.
\section*{Втеча}
Втеча — така сама дія, як інші. Киньте Спритність проти ЦЗ, визначеного Оповідачем: втекти від повільного гобліна — ЦЗ 3, випередити вершника — ЦЗ 5 або більше.
\section*{Протилежні кидки}
Коли два персонажі безпосередньо протидіють одне одному, кидають обидва. Перемагає більший результат. Нічия означає сюжетний глухий кут.
\section*{Допомога іншому протагоністові}
Протагоніст може допомогти союзнику. Ефект залежить від уміння помічника виконувати саме цю дію.

\textbf{Компетентна допомога} (помічник має щонайменше 2d6 у цьому випробуванні): додає +1d6 до набору кубиків союзника.

\textbf{Некомпетентна допомога} (менше ніж 2d6): додає +1d6, але ЦЗ збільшується на 1. Якщо ЦЗ було 6, тепер потрібні дві шістки; якщо вже потрібні дві шістки, тепер потрібні три, і так далі.
\begin{rpgnote}Невмілий помічник зі стародавнім гримуаром може зробити завдання неможливим. Гравець вирішує, чи варто ризикувати.\end{rpgnote}
\section*{Відстані}
Гра не використовує сітку чи точні виміри. Відстані позначають чотирма сюжетними категоріями.
\begin{rpgtable}{Y|Y|Y}
\rpghead{Відстань} & \rpghead{Значення} & \rpghead{Приклади} \\ \hline
Впритул & Фізичний контакт & Боротьба, кинджал, захват \\ \hline
Поруч & Кілька кроків, можна миттєво дістати & Меч, атака, розмова \\ \hline
У межах досяжності & Відстань для стрільби, не для ближнього бою & Лук, дистанційна магія \\ \hline
Далеко & Поза досяжністю; треба наблизитися & За межами пострілу, втеча \\ \hline
\end{rpgtable}
Переміщення на одну категорію ближче — це дія. Оповідач може вимагати кидок, якщо персонаж наближається під тиском.
\section*{Приклад бою}
\begin{rpgnote}
Загін — Кара (Боєць), Сіл (Авантюрист) і Венн (Містик) — потрапляє в засідку двох орків-рейдерів (Експерти; ЦЗ 5 для влучання).

\textbf{Оповідач:} «Із дерев вистрибують два орки. Вони Поруч. У них бойові молоти. Що робите?»

\textbf{Кара:} «Кидаюся на першого з молотом». Сила Експерт (2d6↑): 3 і 5, бере 5. Успіх проти ЦЗ 5. Молот завдає 3 Рани, орк падає з 4 до 1.

\textbf{Сіл:} «Відступаю на дистанцію пострілу й стріляю». Спритність Експерт + Влучність (3d6↑): 2, 4, 6; бере 6. Успіх! Лук завдає 2 Рани, перший орк падає.

\textbf{Венн:} «Стіна полум'я, щоб зупинити другого!» Витрачає 2 Втоми (з 3 до 1 ОЗ, важкий стан). Диво Експерт + Спрямування (3d6↑), але стан дає штраф два щаблі (до 1d6). Випадає 4 проти ЦЗ 5 — невдача. Віддача: полум'я вибухає, усі отримують 2 Рани. Венн падає до −1 ОЗ і вибуває зі сцени.
\end{rpgnote}
\RPGChapter{Зброя та захист}
Шкода залежить від поєднання зброї та Характеристики. Кожна зброя природно відповідає певній Характеристиці; неправильне використання зменшує шкоду.
\section*{Таблиця зброї}
\begin{rpgtable}{Y|Y|Y|Y}
\rpghead{Зброя} & \rpghead{Сила} & \rpghead{Спритність} & \rpghead{Тип} \\ \hline
Кинджал & 1Р & 2Р & Рани \\ \hline
Меч & 2Р & 2Р & Рани \\ \hline
Бойовий молот & 3Р & 1Р & Рани \\ \hline
Лук & — & 2Р & Рани \\ \hline
Спис & 2Р & 2Р & Рани \\ \hline
Бойова сокира & 2Р & 1Р & Рани \\ \hline
Булава & 2В & 1В & Втома \\ \hline
Посох & 1В & 1В & Втома \\ \hline
Кулаки (Боротьба) & 1В & 1В & Втома \\ \hline
\end{rpgtable}
Р = Рана (смертельна шкода). В = Втома (несмертельне ушкодження, оглушення). Персонаж, що опустився до нуля ОЗ лише через Втому, непритомніє, але не ризикує померти.

Критичний успіх у бою завдає додатково \textbf{+1 Рану} понад звичайну шкоду зброї.
\section*{Захист}
Обладунок визначає ЦЗ, якого ворог має досягти, щоб влучити. Кожен обладунок має два значення — проти атак Силою та Спритністю.
\begin{rpgtable}{Y|Y|Y|Y}
\rpghead{Обладунок} & \rpghead{ЦЗ проти Сили} & \rpghead{ЦЗ проти Спритності} & \rpghead{Доступ} \\ \hline
Без обладунку & 3 & 3 & Усі \\ \hline
Легкий & 3 & 4 & Усі класи \\ \hline
Важкий & 4 & 5 & Лише Боєць \\ \hline
\end{rpgtable}
\textbf{Обмеження:} лише Боєць навчений носити важкі обладунки. Авантюристи, Учені та Містики не можуть їх використовувати незалежно від Характеристик.

\textbf{Щит:} коли ворог атакує вас, ЦЗ його атаки підвищується на 1. Потрібні Вміння Щити й вільна рука.
\section*{Визначні предмети}
Деякі предмети мають ім'я та особливий ефект: покращують кидок на один щабель у конкретній ситуації або дозволяють зробити те, що інакше неможливе.
\begin{rpgnote}
\textbf{Соляний клинок} — +1 щабель проти неприродних створінь.
\textbf{Ключ потойбіччя} — відкриває будь-який замок, але кожне використання привертає небажану увагу.
\end{rpgnote}
Це лише приклади. Оповідач створює визначні предмети за тим самим принципом: промовиста назва, конкретний ефект і, можливо, ціна чи наслідок застосування.
\RPGEngravedPanel{Інші визначні предмети}{assets/panel-12-engraved.png}{\textbf{Корона наказу} — +1 щабель до Риторики; власник не може брехати.\par\smallskip \textbf{Ліхтар мертвих} — відкриває невидимих духів і примар, але приваблює ворожих духів.\par\smallskip \textbf{Ковальські рукавиці} — +1 щабель до Руйнування; перегріваються після кожного використання: 1В.}
\RPGChapter{Рани та відновлення}
Очки здоров'я утворюють єдиний запас для бойової шкоди та Ран Втоми — ціни будь-якого надзвичайного зусилля: закляття, піднімання брили чи багатогодинного бігу.
\section*{Два типи ран, один запас}
\textbf{Бойові Рани:} за нуля ОЗ персонаж вибуває зі сцени — непритомніє, може потрапити в полон чи опинитися при смерті. Смерть визначається спільним сюжетним вибором.

\textbf{Рани Втоми:} якщо ОЗ падає до нуля лише через Втому, персонаж непритомніє, але його життю ніщо не загрожує. Він прокидається після короткого відпочинку. Втома і бойові Рани додаються: Містик, який витратив 2 ОЗ на закляття й отримав 1 бойову Рану, справді в небезпеці.
\section*{Стан персонажа}
\begin{rpgtable}{Y|Y|Y}
\rpghead{Залишок ОЗ} & \rpghead{Стан} & \rpghead{Ефект} \\ \hline
Усі & Здоровий & Немає \\ \hline
≤2 & Побитий & −1 щабель до всіх кидків \\ \hline
≤1 & Важкий & −2 щаблі \\ \hline
0 & Поза сценою & Непритомний або при смерті \\ \hline
\end{rpgtable}
\section*{Смерть}
Персонаж помирає, лише якщо \textbf{самі бойові Рани (Р)} дорівнюють або перевищують максимальні ОЗ. Якщо хоча б один пункт шкоди становить Втома (В), персонаж непритомний, але живий.

Лежачого персонажа зі змішаною шкодою можна добити: подальші бойові Рани, що доведуть сумарні Р до максимальних ОЗ, спричинять смерть. Загін має час втрутитися, але смертельний удар можливий.
\section*{Стани}
Отрута, вогонь, сліпота, страх та інші стани не мають окремих правил. Оповідач трактує їх як ситуативний штраф (−1 щабель), доки стан не мине. Та сама механіка охоплює все.
\section*{Відновлення}
\textbf{Короткий відпочинок} (кілька годин у безпеці): відновити 1 Рану.

\textbf{Довгий відпочинок} (ціла ніч у безпеці): відновити всі Рани та скинути правило виснаження.

\textbf{Лікування у сцені:} Вміння Медицина + успіх = відновити 1 Рану (2 у разі Критичного успіху).
\RPGEngravedPanel{Спіраль падіння}{assets/panel-13-engraved.png}{Кара має 6 ОЗ. Орк завдає їй 2Р. Потім вона відчиняє важку браму — 1В. Лишається 3 ОЗ: Побита, −1 щабель до всього. Орк знову б'є — 2Р. Лишається 1 ОЗ: Важкий стан, −2 щаблі. Вибір: утекти чи оголосити Напруження і зробити останній відчайдушний удар. Останній кидок.}
\RPGChapter{Магія та гримуари}
Гримуари — це окремі закляття, які Містик купує під час створення. Кожен працює як картка з визначеними властивостями. Магія не безкоштовна: кожне чаклування коштує Ран Втоми, а невдача викликає \textbf{Віддачу}, яка може вразити Містика, союзників чи довкілля.
\section*{Структура гримуара}
\textbf{Назва} — ім'я закляття.

\textbf{Потрібне Вміння} — Спрямування, Ритуали, Алхімія, Окультизм або Магічне чуття.

\textbf{ЦЗ} — складність чаклування (3–6 або Критичний успіх).

\textbf{Ціна} — Рани Втоми, які сплачують під час чаклування. Легкі гримуари можуть коштувати 0: єдиний ризик — Віддача.

\textbf{Ефект} — наслідок успіху.

\textbf{Віддача} — конкретний наслідок невдачі.
\section*{Приклади гримуарів}
\textbf{Стіна полум'я} — Спрямування, ЦЗ 5, ціна 2.
\textit{Ефект:} стіна вогню перекриває прохід або оточує ворогів.\par
\textit{Віддача:} полум'я виходить з-під контролю. Усі присутні отримують 2 Рани.
\medskip
\textbf{Цілющий дотик} — Спрямування, ЦЗ 3, ціна 1.
\textit{Ефект:} зцілює союзнику 2 Рани.\par
\textit{Віддача:} рани переходять до Містика.
\medskip
\textbf{Коло захисту} — Ритуали, ЦЗ 4, ціна 1.
\textit{Ефект:} обмежена ділянка стає недоступною ворогам на одну сцену.\par
\textit{Віддача:} коло замикає й союзників. Ніхто не може вийти.
\medskip
\textbf{Алхімічний вогонь} — Алхімія, ЦЗ 4, ціна 1.
\textit{Ефект:} флакон полум'я вибухає при зіткненні, завдаючи 2 Рани в області.\par
\textit{Віддача:} вибухає в руках. Ви та всі Поруч отримуєте 1 Рану.
\medskip
\textbf{Голос мертвих} — Окультизм, ЦЗ 5, ціна 2.
\textit{Ефект:} допитуєте духа; він чесно відповідає на три запитання.\par
\textit{Віддача:} ворожий дух тимчасово вселяється в члена загону.
\medskip
\textbf{Око порожнечі} — Магічне чуття, ЦЗ 4, ціна 1.
\textit{Ефект:} на одну сцену бачите крізь стіни й ілюзії.\par
\textit{Віддача:} щось із того боку бачить вас. Воно вас знайшло.
\section*{Створення нових гримуарів}
Оповідач і гравці можуть створювати гримуари за структурою з шести полів. Єдине правило: Віддача мусить бути конкретною і співмірною до Ефекту. Що потужніше закляття, то небезпечніша невдача — не лише для Містика.
\RPGEngravedPanel{Баланс гримуара}{assets/panel-15-engraved.png}{\textbf{ЦЗ 3, ціна 1:} скромний ефект, особиста Віддача. Містик ризикує тільки собою. Добре для щоденних потреб.\par\smallskip\textbf{ЦЗ 5, ціна 2:} потужний ефект, Віддача на весь загін. Ризикують усі; група вирішує разом.\par\smallskip\textbf{Критичний, ціна 3:} диво. Невдача — катастрофа. Використовуйте лише за повної безвиході.}
\section*{Магія для інших класів}
Будь-який персонаж може купувати Гримуари, якщо має щонайменше Компетентний рівень Дива та потрібне Вміння. Проте Диво для всіх, окрім Містика, коштує вдвічі дорожче, тому магія можлива, але дуже дорога.
\RPGChapter{Розвиток}
Немає очок досвіду. Наприкінці сесії загін відповідає на три запитання:
\begin{enumerate}[leftmargin=*,itemsep=2pt]
\item Чи відкрили ви щось важливе?
\item Чи завдали ваша Риса або Тягар вам проблем?
\item Чи зробили ви щось, що змінило історію?
\end{enumerate}
Кожне «так» дає 1 очко розвитку. Їх витрачають так само, як очки створення: на Характеристики, Вміння чи Гримуари.
\section*{Оголошення та розвиток у грі}
Отримавши очко, гравець негайно оголошує, куди його вкладає. Якщо покращення коштує кілька очок, їх вкладають по одному. Тим часом потрібно \textbf{відігравати} розвиток: тренуватися, шукати наставника, експериментувати.

Оповідач може відмовити у розвитку, який жодним чином не відображався під час сесій. Між окремими пригодами з місяцями ігрового часу достатнім поясненням може бути сам час.

Окрім механіки, гравець може:
\begin{itemize}[leftmargin=*,itemsep=2pt]
\item Переписати Рису — герой змінився.
\item Подолати Тягар і вибрати новий.
\end{itemize}
\RPGEngravedPanel{Зростання через гру}{assets/panel-16-engraved.png}{\textbf{Сесія 3:} Сіл вирішує довести Спритність до Виняткового рівня. Це коштує 4 очки; вона вкладає перше.\par\smallskip\textbf{Сесія 5:} знаходить наставника-лучника в лісі й вкладає друге.\par\smallskip\textbf{Сесія 8:} щовечора тренується в таборі — третє.\par\smallskip\textbf{Сесія 11:} проходить неможливе випробування — четверте. Тепер Спритність Виняткова. Вона це прожила.}
\RPGChapter{Вороги}
Ворог не потребує повного аркуша персонажа. У нього лише чотири характеристики:

\textbf{ЦЗ} — два значення (Сила / Спритність), що визначають складність влучання, а також його рівень кидка під час атаки.

\textbf{Рани} — запас здоров'я (1–3 у дрібних ворогів, 4–6 у ветеранів, 7+ у босів).

\textbf{Зброя} — визначає шкоду за тією самою таблицею, що й для героїв.

\textbf{Поведінка} — коротка фраза, яка спрямовує дії ворога.
\section*{Приклади}
\begin{rpgnote}
\textbf{Гоблін} — ЦЗ: С3/Спр3, Компетентний, 2 Рани, кинджал (2Р від Спритності). «Нападає зграєю, тікає наодинці».

\textbf{Орк-рейдер} — ЦЗ: С3/Спр4, Експерт, 4 Рани, бойовий молот (3Р від Сили), легкий обладунок. «Кидається на найслабшого».

\textbf{Зіпсований лицар} — ЦЗ: С4/Спр5, Експерт, 6 Ран, меч (2Р), важкий обладунок. «Шукає двобою, ігнорує слабких».
\end{rpgnote}
\section*{Групи}
Однакові вороги діють як єдиний блок. Кидають один раз, але завдають +1 Рану за кожного учасника понад першого (максимум +3). Кожна Рана, завдана групі, усуває одного її члена.
\section*{Особливі здібності}
Поведінка ворога може містити особливі здібності, сформульовані в одному реченні. Оповідач трактує їх як сюжетні умови.
\begin{rpgnote}
\textbf{Болотяний троль} — ЦЗ: С4/Спр3, Експерт, 6 Ран, кігті (2Р від Сили). «Відновлює 1 Рану щохідно; вогонь зупиняє регенерацію».

\textbf{Виверна} — ЦЗ: С4/Спр4, Експерт, 5 Ран, укус (3Р від Сили). «Літає. Отрута завдає додатково 1 Рану Втоми за кожне влучання».
\end{rpgnote}
\RPGChapter{Напруження}
\begin{rpgnote}«Я стікаю кров'ю, але зроблю це».\end{rpgnote}
Перед кидком гравець може оголосити \textbf{Напруження}: додати \textbf{+1d6} до кидка. \textbf{Після} кидка він отримує \textbf{1 Рану Втоми}. Додатковий кубик впливає і на результат, і на Критичний успіх.

Напруження обмежене +1d6 за дію — ефекти не складаються. Оскільки Рана нараховується після кидка, штрафи станів (Побитий, Важкий) не впливають на поточну дію, лише на наступні.
\begin{rpgnote}Містик із 3 ОЗ витратив 1 на Гримуар, залишилося 2 ОЗ. Для дії потрібне ЦЗ 5. Він оголошує Напруження: додає 1d6 і кидає з повною силою. Після кидка платить 1В, падає до 1 ОЗ (Важкий стан). Якщо провалиться й отримає удар — вибуває зі сцени. Останній кидок.\end{rpgnote}
\RPGEngravedPanel{Напруження поза боєм}{assets/panel-17-engraved.png}{Учений мусить перекласти стародавній напис, перш ніж храм завалиться. Знання Експерт, ЦЗ 5. Кидає 2d6↑: 3 і 4 — невдача.\par\smallskip Оголошує Напруження: 1В, додає 1d6. Тепер три кубики. Випадає 2, 4, 6 — успіх! Але лишається 2 ОЗ: Побитий. Коли підлога обвалиться, доведеться бігти на тремтячих ногах.}
\RPGChapter{Поради оповідачеві}
\begingroup\setlength{\parskip}{5pt plus .5pt minus .5pt}
\textbf{Цільове значення — ваш головний інструмент.} Базове ЦЗ — 4. Підвищуйте його для складних дій, знижуйте для простих. Більше нічого не потрібно.

\textbf{Виснаження створює пригоду.} Провалили дію, що потребує Критичного успіху? Не можна просто повторити — треба шукати інший шлях. Це змушує мислити творчо.

\textbf{Кожен ворог чогось прагне.} Бандит хоче золота, а не смерті; демон — угоди, а не бійки. Мотиви керують конфліктом.

\textbf{Завжди запитуйте: «Як саме?»} Відповідь визначає Характеристику та ЦЗ.

\textbf{Враховуйте ступінь успіху.} 6 проти ЦЗ 3 відрізняється від 4 проти ЦЗ 4. Обидва кидки успішні, але перший — блискучий, другий — ледь-ледь.

\textbf{Рани Втоми — інструмент.} Будь-яке надзвичайне зусилля може коштувати Втоми: підняття валунів, багатогодинний біг, опір отруті. Це не лише для магії.

\textbf{Магія — рішення всієї групи.} Віддача здатна вразити всіх. Нехай гравці обговорюють ризиковані закляття.

\textbf{Соціальні взаємодії підкоряються тим самим правилам.} Переконання, брехня, залякування: Знання + відповідне Вміння проти ЦЗ. Окремих правил не потрібно. Соціальний конфлікт працює як фізичний.
\par\endgroup
\RPGBackCoverUkrainian{ОСТАННІЙ КИДОК}
\end{document}
